\scanpage{65}
$X$ étant projectif et lisse sur $k$ alg. clos, % 1
morphisme structural $f:X\longrightarrow\operatorname{Spec} k$. % 2
Une autre façon est celle-ci % 3
\[
 \boxed{A^i(X)\deleted{\otimes_{\mathbb{Z}}\mathbb{Q}}\simeq
 \operatorname{Im}\bigl[H^{2i}(X,\mathbb{Q}(i))\longrightarrow
 H^{2i}(X,\mathbb{Q}_{\ell}(i))\bigr]} % 4
\]

\uline{NB.} Si $k$ est la clôture algébrique % 5
de $\mathbb{F}_p$, on \unsure{trouve} pour $X$ projectif lisse sur % 6
$k$ % 7
\[
\begin{aligned}
H^{\deleted{\unsure{2i}}i}(X,\overset{\inserted{\nearrow\,\text{\uuline{pas} }\mathbb{Q}_{\ell}(i)!}}{\mathbb{Q}}(j))
={}&\left\{\begin{aligned}
&0 &&\text{si }\unsure{i|j}\ne 2j,2j+1\\ % 8
&A^{\unsure{i|j}}(X)\otimes_{\mathbb{Z}}\mathbb{Q}
 &&\text{si }i=2j\\ % 9
& &&\text{(canonique)}\\ % 10
&A^{\unsure{j|i}}(X)\otimes_{\mathbb{Z}}\mathbb{Q}
 &&\text{si }i=j+1\\ % 11
& &&\text{(\uline{pas} canonique)}
\end{aligned}\right. % 12
\end{aligned}
\]
De plus, pour $i\ne 2j$, l'image de % 13
$H^i(X,\mathbb{Q}(j))$ dans $H^i(X,\mathbb{Q}_{\ell}(j))$ est nulle. % 14

C'est \uline{un fait général} : si $k$ alg. % 15
clos, alors (si $X$ projectif lisse sur $k$) l'image % 16
de $H^{\deleted{\unsure{2i}}j}(X,\mathbb{Q}(\deleted{\unsure{j}}i))\to H^j(X,\mathbb{Q}_{\ell}(i))$ % 17
est \uline{nulle} pour $j\ne 2i$ (question de % 18
poids) ; pour $j=2i$ ce sont les cycles % 19
algébriques \inserted{à coeff} rationnels. % 20
% archive sheet number: 64
